\documentclass[12pt,a4paper]{article}

\usepackage[T2A]{fontenc}
\usepackage[utf8]{inputenc}
\usepackage[russian]{babel}
\usepackage{amsmath,amssymb,amsthm,mathtools}
\usepackage{enumitem}
\usepackage{array}
\usepackage{booktabs}
\usepackage{algorithm}
\usepackage{algorithmic}
\usepackage{geometry}
\geometry{margin=2.2cm}

\newtheorem{definition}{Определение}[section]
\newtheorem{theorem}{Теорема}[section]
\newtheorem{lemma}{Лемма}[section]
\newtheorem{proposition}{Утверждение}[section]
\newtheorem{example}{Пример}[section]
\newtheorem{remark}{Замечание}[section]

\newcommand{\F}{\mathbb{F}}
\newcommand{\E}{\mathbb{E}}
\newcommand{\wt}{ooperatorname{wt}}
\newcommand{\dist}{\operatorname{d}}
\newcommand{\rank}{\operatorname{rank}}
\newcommand{\Ker}{\operatorname{Ker}}
\newcommand{\Imv}{\operatorname{Im}}

\title{Билет №15 \\ \small Коды с исправлением ошибок. Алфавитное кодирование. Методы сжатия информации.  (СПбГУ)}
\author{Сериков Владислав Александрович}
\date{5 мая 2026}


\begin{document}

\maketitle
\tableofcontents
\newpage

\section{Общая схема передачи и хранения информации}

Дискретная передача информации обычно описывается схемой
\[
    \text{источник}
    \longrightarrow
    \text{кодер источника}
    \longrightarrow
    \text{кодер канала}
    \longrightarrow
    \text{канал}
    \longrightarrow
    \text{декодер канала}
    \longrightarrow
    \text{декодер источника}
    \longrightarrow
    \text{получатель}.
\]

Кодер источника устраняет избыточность сообщения, то есть выполняет сжатие.
Кодер канала, наоборот, добавляет специально организованную избыточность, позволяющую обнаруживать и исправлять ошибки.

Различают две основные задачи кодирования:
\begin{enumerate}
    \item \textbf{кодирование источника} --- компактное представление сообщений;
    \item \textbf{канальное кодирование} --- защита сообщений от ошибок при передаче или хранении.
\end{enumerate}

\section{Алфавитное кодирование}

\subsection{Алфавиты, слова и коды}

\begin{definition}
\textbf{Алфавитом} называется конечное непустое множество символов. Обычно исходный алфавит обозначают
\[
    A = \{a_1,\dots,a_m\},
\]
а кодовый алфавит ---
\[
    B = \{b_1,\dots,b_D\}.
\]
Число $D = |B|$ называется основанием кодового алфавита.
\end{definition}

\begin{definition}
Словом над алфавитом $B$ называется конечная последовательность символов из $B$.
Множество всех конечных слов над $B$ обозначается $B^*$.
Длина слова $w$ обозначается $|w|$.
\end{definition}

\begin{definition}
\textbf{Алфавитным кодом} называется отображение
\[
    C : A \to B^*,
\]
которое каждому исходному символу $a_i$ ставит в соответствие кодовое слово
\[
    C(a_i) = c_i \in B^*.
\]
\end{definition}

Отображение $C$ естественным образом продолжается на слова над $A$:
если
\[
    x = a_{i_1}a_{i_2}\dots a_{i_n},
\]
то
\[
    C(x) = C(a_{i_1})C(a_{i_2})\dots C(a_{i_n}).
\]

\begin{definition}
Код называется \textbf{равномерным}, если все кодовые слова имеют одинаковую длину:
\[
    |c_1|=\dots=|c_m|.
\]
Иначе код называется \textbf{неравномерным}.
\end{definition}

\begin{example}
Пусть
\[
    A=\{a,b,c,d\}, \qquad B=\{0,1\}.
\]
Код
\[
    a\mapsto 00,\quad b\mapsto 01,\quad c\mapsto 10,\quad d\mapsto 11
\]
является равномерным двоичным кодом длины $2$.

Код
\[
    a\mapsto 0,\quad b\mapsto 10,\quad c\mapsto 110,\quad d\mapsto 111
\]
является неравномерным.
\end{example}

\subsection{Однозначная декодируемость}

\begin{definition}
Алфавитный код $C:A\to B^*$ называется \textbf{однозначно декодируемым}, если продолжение $C:A^*\to B^*$ является инъективным, то есть из
\[
    C(x)=C(y)
\]
следует
\[
    x=y.
\]
\end{definition}

Иными словами, любую закодированную строку можно разбить на кодовые слова единственным способом.

\begin{example}
Код
\[
    a\mapsto 0,\quad b\mapsto 01
\]
не является однозначно декодируемым, поскольку
\[
    C(ab)=0\,01=001,
\]
а при наличии других слов могут возникать неоднозначности. Более очевидный пример:
\[
    a\mapsto 0,\quad b\mapsto 01,\quad c\mapsto 1.
\]
Тогда
\[
    C(ac)=01=C(b),
\]
следовательно, код не является однозначно декодируемым.
\end{example}

\subsection{Префиксные коды}

\begin{definition}
Слово $u$ называется \textbf{префиксом} слова $v$, если существует слово $w$ такое, что
\[
    v=uw.
\]
\end{definition}

\begin{definition}
Код называется \textbf{префиксным}, или мгновенно декодируемым, если никакое кодовое слово не является префиксом другого кодового слова.
\end{definition}

\begin{example}
Код
\[
    a\mapsto 0,\quad b\mapsto 10,\quad c\mapsto 110,\quad d\mapsto 111
\]
является префиксным.

Код
\[
    a\mapsto 0,\quad b\mapsto 01,\quad c\mapsto 11
\]
не является префиксным, так как слово $0$ является префиксом слова $01$.
\end{example}

\begin{theorem}
Всякий префиксный код является однозначно декодируемым.
\end{theorem}

\begin{proof}
Пусть $C$ --- префиксный код. Предположим, что существует закодированное слово $z$, допускающее два различных разбиения на кодовые слова:
\[
    z=c_{i_1}c_{i_2}\dots c_{i_r}
     =c_{j_1}c_{j_2}\dots c_{j_s}.
\]
Если $c_{i_1}=c_{j_1}$, то можно отбросить общий первый блок и рассматривать оставшуюся часть.
Если первые блоки различны, то оба они начинаются с первого символа строки $z$.
Следовательно, одно из слов $c_{i_1}$ и $c_{j_1}$ должно быть префиксом другого, что противоречит префиксности кода.

Повторяя этот аргумент, получаем, что все блоки в двух разбиениях совпадают, то есть разбиение единственно.
Следовательно, код однозначно декодируем.
\end{proof}

\subsection{Дерево префиксного кода}

Префиксные коды удобно изображать с помощью $D$-арного дерева.
Каждому ребру соответствует символ кодового алфавита.
Кодовыми словам соответствуют листья дерева.
Условие префиксности означает, что кодовые слова не могут соответствовать внутренним вершинам, имеющим потомков.

\begin{example}
Код
\[
    a\mapsto 0,\quad b\mapsto 10,\quad c\mapsto 110,\quad d\mapsto 111
\]
соответствует дереву:
\[
\begin{array}{c}
\text{корень}\\[2mm]
\begin{array}{cc}
0:a & 1\\
    & \begin{array}{cc}
      0:b & 1\\
          & \begin{array}{cc}
            0:c & 1:d
          \end{array}
    \end{array}
\end{array}
\end{array}
\]
\end{example}

\section{Алгоритм Сардинаса--Паттерсона}

\subsection{Назначение алгоритма}

Префиксный код всегда является однозначно декодируемым.
Однако не всякий однозначно декодируемый код является префиксным.

Например, код
\[
    a\mapsto 0,\qquad
    b\mapsto 01,\qquad
    c\mapsto 011
\]
не является префиксным, потому что слово $0$ является префиксом слов $01$ и $011$.
Но вопрос об однозначной декодируемости такого кода требует отдельной проверки.

Алгоритм Сардинаса--Паттерсона позволяет за конечное число шагов определить, является ли заданный конечный алфавитный код однозначно декодируемым.

\subsection{Остатки слов}

Пусть $X$ --- множество кодовых слов:
\[
    X=\{x_1,\dots,x_m\}\subset B^*.
\]

Если слово $u$ является префиксом слова $v$, то существует слово $w$ такое, что
\[
    v=uw.
\]
Слово $w$ называется остатком при удалении префикса $u$ из $v$.

Будем использовать обозначение
\[
    u^{-1}v=w,
\]
если
\[
    v=uw.
\]
Если $u$ не является префиксом $v$, выражение $u^{-1}v$ не определено.

Для множеств слов $A,B\subseteq B^*$ положим
\[
    A^{-1}B
    =
    \{w\in B^*: \exists a\in A,\exists b\in B,\ b=aw\}.
\]
То есть $A^{-1}B$ --- множество всех непустых и пустых остатков, которые получаются при удалении слова из $A$ как префикса слова из $B$.

\subsection{Построение множеств Сардинаса--Паттерсона}

Пусть задан код
\[
    X=\{x_1,\dots,x_m\}.
\]
Положим
\[
    S_1=X^{-1}X\setminus\{\varepsilon\}.
\]
То есть $S_1$ состоит из всех непустых остатков, которые возникают, когда одно кодовое слово является собственным префиксом другого.

Далее рекурсивно определяем множества
\[
    S_{n+1}
    =
    X^{-1}S_n\ \cup\ S_n^{-1}X,
    \qquad n\ge 1.
\]
Здесь:
\begin{itemize}
    \item $X^{-1}S_n$ состоит из остатков, возникающих при удалении кодового слова как префикса слова из $S_n$;
    \item $S_n^{-1}X$ состоит из остатков, возникающих при удалении слова из $S_n$ как префикса кодового слова.
\end{itemize}

Обычно пустое слово $\varepsilon$ рассматривается отдельно:
если на некотором шаге появляется
\[
    \varepsilon\in S_n,
\]
то код не является однозначно декодируемым.

\subsection{Критерий Сардинаса--Паттерсона}

\begin{theorem}[Критерий Сардинаса--Паттерсона]
Пусть $X$ --- конечный код. Тогда $X$ не является однозначно декодируемым тогда и только тогда, когда в процессе построения множеств
\[
    S_1,S_2,\dots
\]
на некотором шаге выполняется одно из условий:
\[
    \varepsilon\in S_n
\]
или, что эквивалентно,
\[
    S_n\cap X\ne \varnothing.
\]
Если же ни одно из этих условий не возникает, а некоторое множество повторилось:
\[
    S_i=S_j,\qquad i<j,
\]
то код является однозначно декодируемым.
\end{theorem}

\begin{proof}
Сначала докажем необходимость.

Пусть код не является однозначно декодируемым.
Тогда существуют две различные последовательности кодовых слов
\[
    x_{i_1}x_{i_2}\dots x_{i_r}
    =
    x_{j_1}x_{j_2}\dots x_{j_s}.
\]
Выберем такую неоднозначность минимальной суммарной длины.
Первые кодовые слова в двух разложениях не могут быть равны, иначе их можно было бы сократить, получив более короткую неоднозначность.
Значит,
\[
    x_{i_1}\ne x_{j_1}.
\]
Но обе стороны равенства начинаются с одного и того же символа, поэтому одно из слов $x_{i_1}$, $x_{j_1}$ является префиксом другого.
Пусть, например,
\[
    x_{j_1}=x_{i_1}u.
\]
Тогда
\[
    u\in S_1.
\]
После сокращения общего префикса $x_{i_1}$ получаем
\[
    x_{i_2}\dots x_{i_r}
    =
    u x_{j_2}\dots x_{j_s}.
\]

Теперь слово $u$ сравнивается со следующими кодовыми словами в другом разложении.
Если $u$ полностью совпадает с некоторым кодовым словом, возникает
\[
    u\in S_n\cap X,
\]
и это означает неоднозначность декодирования.
Если одно из слов является префиксом другого, возникает новый остаток, принадлежащий следующему множеству $S_{n+1}$.

Продолжая процесс, из конечной неоднозначности обязательно получаем либо пустой остаток, либо остаток, совпадающий с кодовым словом.
Следовательно, если код неоднозначно декодируем, то на некотором шаге
\[
    \varepsilon\in S_n
\]
или
\[
    S_n\cap X\ne\varnothing.
\]

Теперь докажем достаточность.

Пусть на некотором шаге
\[
    \varepsilon\in S_n.
\]
Это означает, что после последовательного удаления префиксов была получена полная компенсация двух различных цепочек кодовых слов.
Следовательно, существует слово, допускающее два разных разложения на кодовые слова.
Код не является однозначно декодируемым.

Пусть теперь
\[
    S_n\cap X\ne\varnothing.
\]
Тогда некоторый остаток $u\in S_n$ сам является кодовым словом:
\[
    u\in X.
\]
Но появление $u$ как остатка означает, что существуют две последовательности кодовых слов, одна из которых после сокращения другой даёт хвост $u$.
Так как $u$ тоже кодовое слово, этот хвост можно завершить кодовым словом, получив две разные декомпозиции одного и того же закодированного сообщения.
Следовательно, код не является однозначно декодируемым.

Наконец, поскольку код конечен, все возможные остатки являются суффиксами кодовых слов или их комбинаций ограниченной длины.
Число различных множеств остатков, которые могут возникнуть, конечно.
Если ни $\varepsilon$, ни пересечение с $X$ не появились, но некоторое множество повторилось,
\[
    S_i=S_j,
\]
то дальнейший процесс будет периодическим и новых остатков не возникнет.
Значит, неоднозначность никогда не появится.
Код является однозначно декодируемым.
\end{proof}

\subsection{Алгоритм Сардинаса--Паттерсона}

\begin{algorithm}[H]
\caption{Проверка однозначной декодируемости по Сардинасу--Паттерсону}
\begin{algorithmic}[1]
\STATE Задать множество кодовых слов $X$.
\STATE Построить
\[
    S_1=X^{-1}X\setminus\{\varepsilon\}.
\]
\STATE Если
\[
    S_1\cap X\ne\varnothing,
\]
то код не является однозначно декодируемым.
\STATE Для $n=1,2,\dots$ строить
\[
    S_{n+1}=X^{-1}S_n\cup S_n^{-1}X.
\]
\STATE Если
\[
    \varepsilon\in S_{n+1}
\]
или
\[
    S_{n+1}\cap X\ne\varnothing,
\]
то код не является однозначно декодируемым.
\STATE Если
\[
    S_{n+1}=\varnothing,
\]
то код является однозначно декодируемым.
\STATE Если множество $S_{n+1}$ уже встречалось раньше, то код является однозначно декодируемым.
\end{algorithmic}
\end{algorithm}

\subsection{Пример 1: код не является однозначно декодируемым}

Рассмотрим код
\[
    X=\{0,01,1\}.
\]
Он не является префиксным, так как
\[
    0
\]
является префиксом слова
\[
    01.
\]

Строим множество $S_1$.
Так как
\[
    01=0\cdot 1,
\]
получаем остаток
\[
    1.
\]
Следовательно,
\[
    S_1=\{1\}.
\]
Но
\[
    1\in X.
\]
Значит,
\[
    S_1\cap X\ne\varnothing.
\]
По критерию Сардинаса--Паттерсона код не является однозначно декодируемым.

Действительно,
\[
    01
\]
можно декодировать двумя способами:
\[
    01 = C(b),
\]
если
\[
    b\mapsto 01,
\]
и
\[
    01=C(a)C(c),
\]
если
\[
    a\mapsto 0,\qquad c\mapsto 1.
\]

\subsection{Пример 2: код не префиксный, но однозначно декодируемый}

Рассмотрим код
\[
    X=\{0,01,011\}.
\]
Код не является префиксным, поскольку
\[
    0
\]
является префиксом слов
\[
    01,\quad 011.
\]

Построим множества Сардинаса--Паттерсона.

Сначала:
\[
    S_1=X^{-1}X\setminus\{\varepsilon\}.
\]
Так как
\[
    01=0\cdot 1,
\]
получаем остаток $1$.
Также
\[
    011=0\cdot 11,
\]
получаем остаток $11$.
Кроме того,
\[
    011=01\cdot 1,
\]
получаем остаток $1$.
Следовательно,
\[
    S_1=\{1,11\}.
\]

Проверим пересечение:
\[
    S_1\cap X=\varnothing,
\]
поскольку ни $1$, ни $11$ не являются кодовыми словами.

Теперь строим
\[
    S_2=X^{-1}S_1\cup S_1^{-1}X.
\]

Сначала найдём $X^{-1}S_1$.
Нужно проверить, можно ли удалить кодовое слово из начала слова $1$ или $11$.
Ни одно из кодовых слов
\[
    0,\quad 01,\quad 011
\]
не является префиксом слов $1$ или $11$.
Значит,
\[
    X^{-1}S_1=\varnothing.
\]

Теперь найдём $S_1^{-1}X$.
Нужно проверить, является ли $1$ или $11$ префиксом какого-либо кодового слова из $X$.
Кодовые слова все начинаются с $0$, поэтому это невозможно.
Следовательно,
\[
    S_1^{-1}X=\varnothing.
\]

Итак,
\[
    S_2=\varnothing.
\]
По алгоритму Сардинаса--Паттерсона код является однозначно декодируемым.

Этот пример показывает, что префиксность является достаточным, но не необходимым условием однозначной декодируемости.

\subsection{Пример 3: более подробная проверка}

Рассмотрим код
\[
    X=\{01,10,011,110\}.
\]

Сначала найдём $S_1$.

Проверяем пары кодовых слов:
\[
    011 = 01\cdot 1,
\]
поэтому остаток:
\[
    1.
\]
Также
\[
    110
\]
не начинается с $10$, и $10$ не является префиксом других слов.
Слово $01$ не является префиксом $10$ или $110$.
Слово $10$ не является префиксом других кодовых слов.
Слово $011$ не является префиксом других кодовых слов.
Слово $110$ не является префиксом других кодовых слов.

Следовательно,
\[
    S_1=\{1\}.
\]

Проверяем:
\[
    S_1\cap X=\varnothing.
\]

Теперь строим
\[
    S_2=X^{-1}S_1\cup S_1^{-1}X.
\]

Множество $X^{-1}S_1$ пусто, потому что ни одно кодовое слово из $X$ не является префиксом слова $1$.

Теперь $S_1^{-1}X$.
Слово $1$ является префиксом кодовых слов
\[
    10,\quad 110.
\]
Имеем:
\[
    10=1\cdot 0,
\]
остаток:
\[
    0.
\]
Также
\[
    110=1\cdot 10,
\]
остаток:
\[
    10.
\]
Значит,
\[
    S_2=\{0,10\}.
\]

Но
\[
    10\in X.
\]
Следовательно,
\[
    S_2\cap X\ne\varnothing.
\]
По критерию Сардинаса--Паттерсона код не является однозначно декодируемым.

Действительно, существует неоднозначность:
\[
    0110 = 011\,10 = 01\,110.
\]
То есть одно и то же закодированное слово допускает два различных разбиения на кодовые слова.


\section{Неравенство Крафта--Макмиллана}

\subsection{Неравенство Крафта для префиксных кодов}

\begin{theorem}[Неравенство Крафта]
Пусть существует префиксный код над $D$-ичным кодовым алфавитом с длинами кодовых слов
\[
    l_1,\dots,l_m.
\]
Тогда
\[
    \sum_{i=1}^{m} D^{-l_i} \le 1.
\]
Обратно, если натуральные числа $l_1,\dots,l_m$ удовлетворяют условию
\[
    \sum_{i=1}^{m} D^{-l_i} \le 1,
\]
то существует префиксный $D$-ичный код с такими длинами кодовых слов.
\end{theorem}

\begin{proof}
Пусть $L=\max_i l_i$.
Рассмотрим полное $D$-арное дерево глубины $L$.
На уровне $l_i$ кодовое слово $c_i$ соответствует вершине.
Так как код префиксный, поддеревья, порождённые этими вершинами, не пересекаются.

Вершина уровня $l_i$ имеет ровно $D^{L-l_i}$ листьев на уровне $L$.
Поскольку соответствующие множества листьев не пересекаются, получаем
\[
    \sum_{i=1}^{m} D^{L-l_i} \le D^L.
\]
Делим на $D^L$:
\[
    \sum_{i=1}^{m} D^{-l_i} \le 1.
\]

Докажем обратное утверждение.
Упорядочим длины:
\[
    l_1\le l_2\le\dots\le l_m.
\]
Будем строить кодовые слова в $D$-арном дереве.
На шаге $i$ выбираем самую левую доступную вершину уровня $l_i$, не являющуюся потомком уже выбранных вершин.
Условие
\[
    \sum_{j=1}^{m}D^{-l_j}\le 1
\]
означает, что суммарный объём занятых поддеревьев не превосходит объёма полного дерева.
Следовательно, на каждом шаге такая вершина существует.
Выбранные вершины не находятся друг у друга в потомках, поэтому полученный код является префиксным.
\end{proof}

\subsection{Неравенство Макмиллана}

\begin{theorem}[Неравенство Макмиллана]
Пусть код над $D$-ичным алфавитом является однозначно декодируемым и имеет длины кодовых слов
\[
    l_1,\dots,l_m.
\]
Тогда
\[
    \sum_{i=1}^{m} D^{-l_i} \le 1.
\]
\end{theorem}

\begin{proof}
Обозначим
\[
    K=\sum_{i=1}^{m}D^{-l_i}.
\]
Рассмотрим $n$-ю степень:
\[
    K^n
    =
    \left(\sum_{i=1}^{m}D^{-l_i}\right)^n
    =
    \sum_{i_1,\dots,i_n}
    D^{-(l_{i_1}+\dots+l_{i_n})}.
\]
Положим
\[
    N_s^{(n)}=
    \#\{(i_1,\dots,i_n):l_{i_1}+\dots+l_{i_n}=s\}.
\]
Тогда
\[
    K^n=\sum_s N_s^{(n)}D^{-s}.
\]

Так как код однозначно декодируем, различные последовательности из $n$ исходных символов дают различные кодовые слова.
Кодовых слов длины $s$ над $D$-ичным алфавитом всего $D^s$, поэтому
\[
    N_s^{(n)}\le D^s.
\]
Следовательно,
\[
    K^n
    \le
    \sum_{s=n l_{\min}}^{n l_{\max}}D^sD^{-s}
    =
    n(l_{\max}-l_{\min})+1.
\]
Пусть $K>1$.
Тогда $K^n$ растёт экспоненциально по $n$, а правая часть растёт линейно по $n$.
При достаточно большом $n$ получаем противоречие.
Следовательно,
\[
    K\le 1.
\]
\end{proof}

\section{Энтропия и средняя длина кода}

\subsection{Энтропия дискретного источника}

Пусть источник выдаёт символы
\[
    a_1,\dots,a_m
\]
с вероятностями
\[
    p_1,\dots,p_m,\qquad p_i>0,\qquad \sum_{i=1}^{m}p_i=1.
\]

\begin{definition}
Энтропией источника называется величина
\[
    H_D(X)
    =
    -\sum_{i=1}^{m}p_i\log_D p_i
    =
    \sum_{i=1}^{m}p_i\log_D\frac1{p_i}.
\]
Если $D=2$, энтропия измеряется в битах.
\end{definition}

\begin{definition}
Средней длиной алфавитного кода называется
\[
    \overline{l}
    =
    \sum_{i=1}^{m}p_i l_i,
\]
где $l_i$ --- длина кодового слова для символа $a_i$.
\end{definition}

\subsection{Нижняя граница средней длины}

\begin{theorem}
Для любого однозначно декодируемого $D$-ичного кода со средней длиной $\overline l$ выполнено
\[
    \overline l \ge H_D(X).
\]
\end{theorem}

\begin{proof}
Пусть длины кодовых слов равны $l_1,\dots,l_m$.
По неравенству Макмиллана
\[
    K=\sum_{i=1}^{m}D^{-l_i}\le 1.
\]
Определим числа
\[
    q_i=\frac{D^{-l_i}}{K}.
\]
Тогда
\[
    q_i>0,\qquad \sum_{i=1}^{m}q_i=1.
\]
Используем неотрицательность относительной энтропии:
\[
    \sum_{i=1}^{m}p_i\log_D\frac{p_i}{q_i}\ge 0.
\]
Отсюда
\[
    \sum_i p_i\log_D p_i
    \ge
    \sum_i p_i\log_D q_i.
\]
Следовательно,
\[
    H_D(X)
    =
    -\sum_i p_i\log_D p_i
    \le
    -\sum_i p_i\log_D q_i.
\]
Но
\[
    \log_D q_i
    =
    \log_D(D^{-l_i})-\log_D K
    =
    -l_i-\log_D K.
\]
Поэтому
\[
    -\sum_i p_i\log_D q_i
    =
    \sum_i p_i l_i+\log_D K.
\]
Так как $K\le 1$, имеем $\log_D K\le 0$.
Значит,
\[
    H_D(X)
    \le
    \overline l+\log_D K
    \le
    \overline l.
\]
Итак,
\[
    \overline l\ge H_D(X).
\]
\end{proof}

\subsection{Теорема Шеннона о кодировании источника}

\begin{theorem}[Шеннон, кодирование источника для одного символа]
Для дискретного источника с энтропией $H_D(X)$ существует префиксный $D$-ичный код со средней длиной $\overline l$, удовлетворяющей
\[
    H_D(X)\le \overline l < H_D(X)+1.
\]
\end{theorem}

\begin{proof}
Нижняя оценка уже доказана для любого однозначно декодируемого кода.

Для верхней оценки положим
\[
    l_i=\left\lceil \log_D\frac1{p_i}\right\rceil.
\]
Тогда
\[
    l_i < \log_D\frac1{p_i}+1.
\]
Следовательно,
\[
    D^{-l_i}\le p_i.
\]
Поэтому
\[
    \sum_{i=1}^{m}D^{-l_i}\le \sum_{i=1}^{m}p_i=1.
\]
По неравенству Крафта существует префиксный код с такими длинами.

Его средняя длина удовлетворяет
\[
    \overline l
    =
    \sum_i p_i l_i
    <
    \sum_i p_i\left(\log_D\frac1{p_i}+1\right)
    =
    H_D(X)+1.
\]
Теорема доказана.
\end{proof}

\section{Эффективность алфавитного кода}

\subsection{Избыточность и эффективность}

Пусть дискретный источник выдаёт символы
\[
    a_1,\dots,a_m
\]
с вероятностями
\[
    p_1,\dots,p_m,
    \qquad
    \sum_{i=1}^{m}p_i=1.
\]
Пусть для этого источника построен $D$-ичный однозначно декодируемый код с длинами кодовых слов
\[
    l_1,\dots,l_m.
\]
Средняя длина кода равна
\[
    \overline l = \sum_{i=1}^{m}p_i l_i.
\]
Энтропия источника по основанию $D$ равна
\[
    H_D(X)=\sum_{i=1}^{m}p_i\log_D\frac1{p_i}.
\]

Из теоремы Шеннона для кодирования источника известно, что для любого однозначно декодируемого кода
\[
    \overline l\ge H_D(X).
\]
То есть энтропия задаёт нижнюю теоретическую границу средней длины кода.

\begin{definition}
\textbf{Эффективностью кода} называется величина
\[
    \eta = \frac{H_D(X)}{\overline l}.
\]
\end{definition}

Так как
\[
    \overline l\ge H_D(X),
\]
то
\[
    0<\eta\le 1.
\]

Если
\[
    \eta=1,
\]
то код достигает энтропийной нижней границы и является идеальным с точки зрения средней длины.

\begin{definition}
\textbf{Абсолютной избыточностью кода} называется величина
\[
    r = \overline l - H_D(X).
\]
\end{definition}

\begin{definition}
\textbf{Относительной избыточностью кода} называется величина
\[
    \rho = 1-\eta
    =
    1-\frac{H_D(X)}{\overline l}
    =
    \frac{\overline l-H_D(X)}{\overline l}.
\]
\end{definition}

Таким образом:
\[
    \eta=\frac{H_D(X)}{\overline l},
    \qquad
    \rho=\frac{\overline l-H_D(X)}{\overline l}.
\]

Эффективность показывает, какая доля средней длины кодового слова теоретически необходима для представления информации, а относительная избыточность показывает долю ``лишних'' символов в среднем кодовом слове.

\begin{example}
Пусть источник имеет вероятности
\[
    p(a)=0.4,\quad p(b)=0.3,\quad p(c)=0.2,\quad p(d)=0.1.
\]
Для кода Хаффмана
\[
    a\mapsto 0,\quad
    b\mapsto 10,\quad
    c\mapsto 110,\quad
    d\mapsto 111
\]
средняя длина равна
\[
    \overline l
    =
    0.4\cdot 1+0.3\cdot 2+0.2\cdot 3+0.1\cdot 3
    =
    1.9.
\]
Энтропия:
\[
\begin{aligned}
H_2(X)
&=
-0.4\log_2 0.4
-0.3\log_2 0.3
-0.2\log_2 0.2
-0.1\log_2 0.1\\
&\approx 1.846.
\end{aligned}
\]
Тогда эффективность кода:
\[
    \eta=\frac{1.846}{1.9}\approx 0.9716.
\]
Относительная избыточность:
\[
    \rho=1-\eta\approx 0.0284.
\]
То есть эффективность кода составляет примерно $97.16\%$, а избыточность --- примерно $2.84\%$.
\end{example}

\subsection{Когда эффективность равна единице}

Эффективность кода равна единице, если
\[
    \overline l = H_D(X).
\]
Для префиксного кода это возможно, например, когда все вероятности символов имеют вид
\[
    p_i=D^{-l_i}
\]
для некоторых натуральных длин $l_i$, удовлетворяющих равенству Крафта
\[
    \sum_{i=1}^{m}D^{-l_i}=1.
\]

\begin{example}
Пусть
\[
    p(a)=\frac12,\qquad
    p(b)=\frac14,\qquad
    p(c)=\frac14.
\]
Возьмём двоичный префиксный код
\[
    a\mapsto 0,\qquad
    b\mapsto 10,\qquad
    c\mapsto 11.
\]
Тогда
\[
    \overline l
    =
    \frac12\cdot 1+\frac14\cdot 2+\frac14\cdot 2
    =
    \frac32.
\]
Энтропия:
\[
\begin{aligned}
H_2(X)
&=
-\frac12\log_2\frac12
-\frac14\log_2\frac14
-\frac14\log_2\frac14\\
&=
\frac12+ \frac12+\frac12
=
\frac32.
\end{aligned}
\]
Следовательно,
\[
    \eta=\frac{H_2(X)}{\overline l}=1.
\]
Код является оптимальным в энтропийном смысле.
\end{example}

\section{Код Хаффмана}

\subsection{Идея метода}

Код Хаффмана строит оптимальный префиксный код для заданных вероятностей символов.
Оптимальность понимается в смысле минимальности средней длины
\[
    \overline l=\sum_i p_i l_i
\]
среди всех префиксных кодов над фиксированным кодовым алфавитом.

Для простоты сначала рассматривается двоичный код.

\subsection{Алгоритм Хаффмана для двоичного алфавита}

\begin{algorithm}[H]
\caption{Построение двоичного кода Хаффмана}
\begin{algorithmic}[1]
\STATE Для каждого символа создать лист дерева с весом, равным вероятности символа.
\STATE Пока в наборе деревьев больше одного дерева:
\begin{enumerate}
    \item выбрать два дерева с наименьшими весами;
    \item объединить их в новое дерево;
    \item вес нового дерева положить равным сумме весов объединённых деревьев;
    \item удалить два выбранных дерева и добавить новое.
\end{enumerate}
\STATE Полученное дерево является кодовым деревом.
\STATE Пометить каждое левое ребро символом $0$, каждое правое ребро символом $1$.
\STATE Кодовое слово символа --- последовательность меток на пути от корня к соответствующему листу.
\end{algorithmic}
\end{algorithm}

\subsection{Минимальный пример}

Пусть алфавит
\[
    A=\{a,b,c,d\}
\]
имеет вероятности
\[
    p(a)=0.4,\quad p(b)=0.3,\quad p(c)=0.2,\quad p(d)=0.1.
\]

Строим дерево Хаффмана.

\[
\begin{array}{c|c}
\text{Шаг} & \text{Действие}\\
\hline
1 & c(0.2) \text{ и } d(0.1) \text{ объединяются в узел } cd(0.3)\\
2 & b(0.3) \text{ и } cd(0.3) \text{ объединяются в узел } bcd(0.6)\\
3 & a(0.4) \text{ и } bcd(0.6) \text{ объединяются в корень }(1.0)
\end{array}
\]

Один из возможных кодов:
\[
    a\mapsto 0,\quad
    b\mapsto 10,\quad
    c\mapsto 110,\quad
    d\mapsto 111.
\]

Средняя длина:
\[
    \overline l
    =
    0.4\cdot 1+0.3\cdot 2+0.2\cdot 3+0.1\cdot 3
    =
    1.9.
\]

Энтропия:
\[
\begin{aligned}
H_2(X)
&=
-0.4\log_2 0.4
-0.3\log_2 0.3
-0.2\log_2 0.2
-0.1\log_2 0.1\\
&\approx 1.846.
\end{aligned}
\]
Действительно,
\[
    H_2(X)\le \overline l < H_2(X)+1.
\]

\subsection{Оптимальность кода Хаффмана}

\begin{lemma}
В некотором оптимальном двоичном префиксном коде два наименее вероятных символа имеют кодовые слова одинаковой максимальной длины и отличаются только последним битом.
\end{lemma}

\begin{proof}
Рассмотрим оптимальное префиксное кодовое дерево.
Если в дереве есть внутренний узел только с одним потомком, его можно удалить, укоротив некоторые кодовые слова и не нарушив префиксность.
Значит, оптимальное дерево можно считать полным: каждый внутренний узел имеет двух потомков.

Возьмём два листа максимальной глубины, являющиеся братьями.
Такие листья существуют в полном двоичном дереве.
Пусть соответствующие им символы имеют вероятности $p_i$ и $p_j$.

Если один из наименее вероятных символов расположен выше, чем лист максимальной глубины, можно поменять его местами с более вероятным символом на максимальной глубине.
Средняя длина при этом не увеличится, поскольку менее вероятный символ получает не меньшую длину, а более вероятный --- не большую.
Следовательно, в некотором оптимальном дереве два наименее вероятных символа находятся в листьях максимальной глубины.

Так как листья максимальной глубины можно выбрать братьями, соответствующие кодовые слова имеют одинаковую длину и отличаются только последним битом.
\end{proof}

\begin{theorem}
Двоичный алгоритм Хаффмана строит префиксный код с минимальной средней длиной среди всех двоичных префиксных кодов для данного распределения вероятностей.
\end{theorem}

\begin{proof}
Докажем по индукции по числу символов $m$.

При $m=2$ оптимальный код очевиден:
\[
    a_1\mapsto 0,\qquad a_2\mapsto 1.
\]

Пусть утверждение верно для алфавитов мощности $m-1$.
Рассмотрим алфавит мощности $m$.
Пусть $x$ и $y$ --- два наименее вероятных символа с вероятностями $p_x$ и $p_y$.
По лемме существует оптимальный код, в котором $x$ и $y$ имеют кодовые слова вида
\[
    u0,\qquad u1.
\]

Заменим символы $x$ и $y$ одним новым символом $z$ с вероятностью
\[
    p_z=p_x+p_y.
\]
Получим задачу для $m-1$ символа.
Если для сокращённого алфавита имеется код со средней длиной $L'$, то после замены $z$ на пару $x,y$ средняя длина увеличится на
\[
    p_x+p_y,
\]
поскольку кодовые слова $x$ и $y$ на один бит длиннее слова $z$.
Итак,
\[
    L = L' + p_x+p_y.
\]

По предположению индукции алгоритм Хаффмана оптимален для сокращённого алфавита.
Следовательно, после обратного разбиения символа $z$ на $x$ и $y$ получаем оптимальный код для исходного алфавита.

Алгоритм Хаффмана именно так и действует: на первом шаге объединяет два наименее вероятных символа, затем рекурсивно строит оптимальное дерево.
Теорема доказана.
\end{proof}

\section{Арифметическое кодирование}

\subsection{Основная идея}

Хаффманов код присваивает целое число бит каждому символу.
Из-за целочисленности длин кодовых слов средняя длина не всегда близка к энтропии.

Арифметическое кодирование кодирует не отдельные символы, а всё сообщение целиком.
Сообщению сопоставляется подынтервал отрезка $[0,1)$.
Чем вероятнее сообщение, тем длиннее соответствующий интервал и тем меньше бит требуется для указания точки внутри него.

\subsection{Алгоритм арифметического кодирования}

Пусть алфавит
\[
    A=\{a_1,\dots,a_m\}
\]
имеет вероятности
\[
    p_1,\dots,p_m.
\]
Сначала строятся накопленные вероятности:
\[
    F_0=0,\qquad
    F_i=\sum_{j=1}^{i}p_j.
\]
Символу $a_i$ соответствует интервал
\[
    [F_{i-1},F_i).
\]

\begin{algorithm}[H]
\caption{Арифметическое кодирование}
\begin{algorithmic}[1]
\STATE Положить $L=0$, $R=1$.
\STATE Для каждого символа $x$ сообщения:
\begin{enumerate}
    \item положить $W=R-L$;
    \item если $x=a_i$, заменить
    \[
        R \leftarrow L + W F_i,
    \]
    \[
        L \leftarrow L + W F_{i-1}.
    \]
\end{enumerate}
\STATE Выбрать любое число $z\in[L,R)$.
\STATE Закодировать двоичную запись числа $z$ с достаточным числом бит.
\end{algorithmic}
\end{algorithm}

\subsection{Пример}

Пусть
\[
    A=\{a,b\},\qquad p(a)=0.75,\quad p(b)=0.25.
\]
Тогда
\[
    a\leftrightarrow [0,0.75),\qquad
    b\leftrightarrow [0.75,1).
\]

Закодируем сообщение
\[
    aba.
\]

Начальный интервал:
\[
    [L,R)=[0,1).
\]

После символа $a$:
\[
    [L,R)=[0,0.75).
\]

После символа $b$ ширина равна $0.75$:
\[
    L=0+0.75\cdot 0.75=0.5625,
\]
\[
    R=0+0.75\cdot 1=0.75.
\]
Получаем
\[
    [0.5625,0.75).
\]

После символа $a$ ширина равна
\[
    0.75-0.5625=0.1875.
\]
Тогда
\[
    L=0.5625+0.1875\cdot 0=0.5625,
\]
\[
    R=0.5625+0.1875\cdot 0.75=0.703125.
\]
Итоговый интервал:
\[
    [0.5625,0.703125).
\]
Можно выбрать, например,
\[
    z=0.625.
\]
В двоичной записи
\[
    0.625=0.101_2.
\]
Значит, одним из возможных кодовых представлений является строка $101$, если известна длина исходного сообщения или используется специальный символ конца сообщения.

\section{Методы сжатия информации}

\subsection{Сжатие без потерь и с потерями}

\begin{definition}
\textbf{Сжатие без потерь} --- это преобразование данных, при котором исходные данные могут быть восстановлены точно.
\end{definition}

Примеры:
\begin{itemize}
    \item кодирование Хаффмана;
    \item арифметическое кодирование;
    \item LZ77, LZ78, LZW;
    \item RLE.
\end{itemize}

\begin{definition}
\textbf{Сжатие с потерями} --- это преобразование данных, при котором восстановленные данные могут отличаться от исходных, но различие считается допустимым.
\end{definition}

Примеры:
\begin{itemize}
    \item JPEG для изображений;
    \item MP3, AAC для звука;
    \item MPEG, H.264, H.265 для видео.
\end{itemize}

\subsection{RLE: кодирование серий}

RLE, или run-length encoding, эффективно сжимает данные, содержащие длинные последовательности одинаковых символов.

\begin{algorithm}[H]
\caption{RLE-кодирование}
\begin{algorithmic}[1]
\STATE Просматривать входную строку слева направо.
\STATE Разбивать строку на максимальные серии одинаковых символов.
\STATE Каждую серию вида
\[
    aaa\dots a
\]
длины $k$ заменять парой
\[
    (a,k).
\]
\end{algorithmic}
\end{algorithm}

\begin{example}
Строка
\[
    AAAAABBBCCDAA
\]
разбивается на серии:
\[
    A^5B^3C^2D^1A^2.
\]
RLE-запись:
\[
    (A,5)(B,3)(C,2)(D,1)(A,2).
\]
\end{example}

RLE хорошо работает для простых изображений, масок, факсов, но плохо работает для случайных данных.

\subsection{LZ77}

Метод LZ77 использует повторения в уже просмотренной части текста.
Кодер поддерживает скользящее окно и заменяет повторяющиеся фрагменты ссылками назад.

Обычно тройка LZ77 имеет вид
\[
    (d,l,c),
\]
где:
\begin{itemize}
    \item $d$ --- расстояние назад до начала совпадения;
    \item $l$ --- длина совпадения;
    \item $c$ --- следующий символ после совпадения.
\end{itemize}

\begin{algorithm}[H]
\caption{Упрощённое LZ77-кодирование}
\begin{algorithmic}[1]
\STATE Установить текущую позицию $i=1$.
\STATE В уже просмотренной части строки найти самый длинный фрагмент, совпадающий с началом строки в позиции $i$.
\STATE Если найдено совпадение длины $l>0$ на расстоянии $d$, вывести тройку $(d,l,c)$, где $c$ --- следующий символ.
\STATE Если совпадение не найдено, вывести $(0,0,x_i)$.
\STATE Сдвинуть позицию на $l+1$ символов.
\STATE Повторять, пока строка не закончится.
\end{algorithmic}
\end{algorithm}

\begin{example}
Рассмотрим строку
\[
    ABABABA.
\]

Упрощённое кодирование:
\[
\begin{array}{c|c|c}
\text{Позиция} & \text{Просмотренная часть} & \text{Код}\\
\hline
1 & \varnothing & (0,0,A)\\
2 & A & (0,0,B)\\
3 & AB & (2,4,A)
\end{array}
\]
На позиции $3$ начало строки равно
\[
    ABABA.
\]
В просмотренной части уже есть фрагмент, начинающийся на расстоянии $2$ назад.
Получается ссылка на повторение.
\end{example}

\subsection{LZ78}

Метод LZ78 строит словарь фраз.
Каждая новая фраза получается добавлением одного символа к уже известной фразе.

Кодовая пара имеет вид
\[
    (j,c),
\]
где $j$ --- номер фразы в словаре, а $c$ --- новый символ.

\begin{algorithm}[H]
\caption{LZ78-кодирование}
\begin{algorithmic}[1]
\STATE Инициализировать словарь пустой строкой с номером $0$.
\STATE Положить текущую фразу $w=\varepsilon$.
\STATE Читать символы слева направо.
\STATE Пока $wc$ уже есть в словаре, заменить $w$ на $wc$ и читать следующий символ $c$.
\STATE Когда $wc$ отсутствует в словаре, вывести пару $(\operatorname{index}(w),c)$.
\STATE Добавить $wc$ в словарь.
\STATE Положить $w=\varepsilon$ и продолжить.
\end{algorithmic}
\end{algorithm}

\begin{example}
Закодируем строку
\[
    ABABABA.
\]

\[
\begin{array}{c|c|c}
\text{Шаг} & \text{Вывод} & \text{Добавление в словарь}\\
\hline
1 & (0,A) & 1:A\\
2 & (0,B) & 2:B\\
3 & (1,B) & 3:AB\\
4 & (3,A) & 4:ABA
\end{array}
\]

Итог:
\[
    (0,A),(0,B),(1,B),(3,A).
\]
\end{example}

\subsection{LZW}

LZW является модификацией LZ78.
В словарь заранее помещаются все односимвольные строки.
На выход передаются только номера словарных фраз.

\begin{algorithm}[H]
\caption{LZW-кодирование}
\begin{algorithmic}[1]
\STATE Инициализировать словарь всеми символами алфавита.
\STATE Прочитать первый символ и положить $w$ равным этому символу.
\STATE Пока вход не закончился:
\begin{enumerate}
    \item прочитать следующий символ $c$;
    \item если $wc$ есть в словаре, положить $w\leftarrow wc$;
    \item иначе вывести код $w$, добавить $wc$ в словарь, положить $w\leftarrow c$.
\end{enumerate}
\STATE В конце вывести код $w$.
\end{algorithmic}
\end{algorithm}

\begin{example}
Пусть начальный словарь:
\[
    A\mapsto 1,\qquad B\mapsto 2.
\]
Кодируем строку
\[
    ABABABA.
\]

\[
\begin{array}{c|c|c}
\text{Шаг} & \text{Вывод} & \text{Добавление}\\
\hline
1 & 1(A) & 3:AB\\
2 & 2(B) & 4:BA\\
3 & 3(AB) & 5:ABA\\
4 & 5(ABA) & -
\end{array}
\]

Итоговая последовательность кодов:
\[
    1,2,3,5.
\]
\end{example}

\section{Коды с обнаружением и исправлением ошибок}

\subsection{Кодовое расстояние}

Пусть рассматриваются слова длины $n$ над алфавитом $Q$.

\begin{definition}
\textbf{Расстоянием Хэмминга} между словами
\[
    x=(x_1,\dots,x_n),\qquad y=(y_1,\dots,y_n)
\]
называется число позиций, в которых они различаются:
\[
    d(x,y)=|\{i:x_i\ne y_i\}|.
\]
\end{definition}

\begin{definition}
Кодом длины $n$ над алфавитом $Q$ называется подмножество
\[
    C\subseteq Q^n.
\]
Элементы $C$ называются кодовыми словами.
\end{definition}

\begin{definition}
Минимальным расстоянием кода $C$ называется
\[
    d(C)=\min_{\substack{x,y\in C\\x\ne y}}d(x,y).
\]
\end{definition}

\begin{definition}
Шаром Хэмминга радиуса $t$ с центром $x$ называется множество
\[
    B_t(x)=\{y\in Q^n:d(x,y)\le t\}.
\]
\end{definition}

\subsection{Обнаружение и исправление ошибок}

\begin{theorem}
Код с минимальным расстоянием $d$ обнаруживает любые $s$ ошибок тогда и только тогда, когда
\[
    d\ge s+1.
\]
\end{theorem}

\begin{proof}
Если при передаче кодового слова произошло не более $s$ ошибок, принятое слово находится на расстоянии не более $s$ от переданного.
Если $d\ge s+1$, то никакое другое кодовое слово не может находиться на расстоянии $\le s$ от переданного.
Следовательно, при наличии ошибки принятое слово не совпадёт с другим кодовым словом, и ошибка будет обнаружена.

Обратно, если $d\le s$, существуют два различных кодовых слова $x,y$ такие, что
\[
    d(x,y)\le s.
\]
Тогда при передаче $x$ можно получить $y$ не более чем за $s$ ошибок.
Поскольку $y$ само является кодовым словом, ошибка не будет обнаружена.
\end{proof}

\begin{theorem}
Код с минимальным расстоянием $d$ исправляет любые $t$ ошибок тогда и только тогда, когда
\[
    d\ge 2t+1.
\]
\end{theorem}

\begin{proof}
Код исправляет $t$ ошибок, если шары Хэмминга радиуса $t$ вокруг различных кодовых слов не пересекаются.

Пусть $d\ge 2t+1$.
Если существует слово
\[
    z\in B_t(x)\cap B_t(y)
\]
для различных кодовых слов $x,y$, то по неравенству треугольника
\[
    d(x,y)\le d(x,z)+d(z,y)\le t+t=2t,
\]
что противоречит $d(x,y)\ge 2t+1$.
Значит, шары не пересекаются, и ближайшее кодовое слово определяется однозначно.

Обратно, если $d\le 2t$, существуют различные кодовые слова $x,y$ такие, что
\[
    d(x,y)\le 2t.
\]
Можно выбрать слово $z$, лежащее на кратчайшем пути между $x$ и $y$, так что
\[
    d(x,z)\le t,\qquad d(y,z)\le t.
\]
Тогда одно и то же принятое слово $z$ могло получиться как из $x$, так и из $y$ при не более чем $t$ ошибках.
Однозначное исправление невозможно.
\end{proof}

\subsection{Параметры кодов}

Код обычно обозначают как
\[
    (n,M,d)_q,
\]
где:
\begin{itemize}
    \item $n$ --- длина кодовых слов;
    \item $M=|C|$ --- число кодовых слов;
    \item $d$ --- минимальное расстояние;
    \item $q=|Q|$ --- мощность алфавита.
\end{itemize}

Для линейных кодов используется обозначение
\[
    [n,k,d]_q,
\]
где
\[
    M=q^k.
\]

\section{Линейные коды}

\subsection{Определение}

Пусть $Q=\F_q$ --- конечное поле из $q$ элементов.

\begin{definition}
Линейным $q$-ичным кодом длины $n$ называется линейное подпространство
\[
    C\le \F_q^n.
\]
Если
\[
    \dim C=k,
\]
то код содержит
\[
    |C|=q^k
\]
слов.
\end{definition}

\begin{definition}
Весом Хэмминга слова $x\in \F_q^n$ называется число его ненулевых координат:
\[
    \wt(x)=|\{i:x_i\ne 0\}|.
\]
\end{definition}

\begin{theorem}
Для линейного кода $C$ минимальное расстояние равно минимальному ненулевому весу:
\[
    d(C)=\min_{\substack{c\in C\\c\ne 0}}\wt(c).
\]
\end{theorem}

\begin{proof}
Для любых $x,y\in C$ имеем
\[
    d(x,y)=\wt(x-y).
\]
Так как $C$ линейный, то
\[
    x-y\in C.
\]
Если $x\ne y$, то $x-y\ne 0$.
Следовательно,
\[
    d(C)
    =
    \min_{\substack{x,y\in C\\x\ne y}}\wt(x-y)
    =
    \min_{\substack{c\in C\\c\ne 0}}\wt(c).
\]
\end{proof}

\subsection{Порождающая матрица}

\begin{definition}
Порождающей матрицей линейного $[n,k]_q$-кода называется матрица
\[
    G\in \F_q^{k\times n},
\]
строки которой образуют базис кода.
\end{definition}

Кодирование сообщения
\[
    u\in \F_q^k
\]
производится по формуле
\[
    c=uG.
\]

Если
\[
    G=(I_k\mid A),
\]
то код называется систематическим: первые $k$ координат кодового слова совпадают с информационными символами.

\subsection{Проверочная матрица}

\begin{definition}
Проверочной матрицей линейного кода $C$ называется матрица
\[
    H\in \F_q^{(n-k)\times n}
\]
такая, что
\[
    C=\{x\in \F_q^n:Hx^T=0\}.
\]
\end{definition}

Если
\[
    G=(I_k\mid A),
\]
то можно взять
\[
    H=(-A^T\mid I_{n-k}).
\]

Действительно,
\[
    GH^T=(I_k\mid A)
    \begin{pmatrix}
    -A\\
    I_{n-k}
    \end{pmatrix}
    =
    -A+A=0.
\]

\subsection{Синдром}

\begin{definition}
Пусть передано кодовое слово $c\in C$, а принято
\[
    y=c+e,
\]
где $e$ --- вектор ошибки.
Синдромом принятого слова называется
\[
    s=Hy^T.
\]
\end{definition}

Так как
\[
    Hc^T=0,
\]
получаем
\[
    s=H(c+e)^T=Hc^T+He^T=He^T.
\]
Синдром зависит только от ошибки, а не от переданного кодового слова.

\begin{algorithm}[H]
\caption{Декодирование по синдрому}
\begin{algorithmic}[1]
\STATE Для каждого возможного исправляемого вектора ошибки $e$ заранее вычислить синдром $He^T$.
\STATE При получении слова $y$ вычислить $s=Hy^T$.
\STATE Найти в таблице ошибку $e$ с синдромом $s$.
\STATE Исправить слово:
\[
    c=y-e.
\]
\end{algorithmic}
\end{algorithm}

\section{Код Хэмминга}

\subsection{Построение двоичного кода Хэмминга}

Для числа проверочных символов $r$ двоичный код Хэмминга имеет параметры
\[
    [2^r-1,\;2^r-r-1,\;3]_2.
\]
Его проверочная матрица $H$ имеет размер
\[
    r\times (2^r-1)
\]
и состоит из всех ненулевых двоичных столбцов длины $r$.

\subsection{Пример: код Хэмминга \texorpdfstring{$[7,4,3]$}{[7,4,3]}}

Возьмём
\[
H=
\begin{pmatrix}
1&0&1&0&1&0&1\\
0&1&1&0&0&1&1\\
0&0&0&1&1&1&1
\end{pmatrix}.
\]
Столбцы матрицы $H$ являются двоичными записями чисел от $1$ до $7$.

Код:
\[
    C=\{x\in \F_2^7:Hx^T=0\}.
\]

Если произошла одиночная ошибка в позиции $i$, то
\[
    e=(0,\dots,0,1,0,\dots,0),
\]
и синдром равен $i$-му столбцу матрицы $H$:
\[
    s=He^T=H_i.
\]
Следовательно, синдром указывает номер ошибочной позиции.

\subsection{Минимальный пример декодирования}

Пусть принято слово
\[
    y=1011010.
\]
Вычисляем синдром:
\[
s=Hy^T.
\]
Имеем
\[
y^T=
\begin{pmatrix}
1\\0\\1\\1\\0\\1\\0
\end{pmatrix}.
\]
Тогда
\[
s=
\begin{pmatrix}
1&0&1&0&1&0&1\\
0&1&1&0&0&1&1\\
0&0&0&1&1&1&1
\end{pmatrix}
\begin{pmatrix}
1\\0\\1\\1\\0\\1\\0
\end{pmatrix}
=
\begin{pmatrix}
1+1\\
1+1\\
1+1
\end{pmatrix}
=
\begin{pmatrix}
0\\0\\0
\end{pmatrix}
\pmod 2.
\]
Значит, ошибок не обнаружено.

Если вместо этого принято
\[
    y'=1011110,
\]
то
\[
y'^T=
\begin{pmatrix}
1\\0\\1\\1\\1\\1\\0
\end{pmatrix}.
\]
Синдром:
\[
s'=Hy'^T
=
\begin{pmatrix}
1+1+1\\
1+1\\
1+1+1
\end{pmatrix}
=
\begin{pmatrix}
1\\0\\1
\end{pmatrix}.
\]
Это пятый столбец матрицы $H$.
Следовательно, ошибка в позиции $5$.
Исправленное слово:
\[
    1011110 + 0000100 = 1011010.
\]

\subsection{Минимальное расстояние кода Хэмминга}

\begin{theorem}
Двоичный код Хэмминга с проверочной матрицей, состоящей из всех ненулевых столбцов длины $r$, имеет минимальное расстояние
\[
    d=3.
\]
\end{theorem}

\begin{proof}
Минимальное расстояние линейного кода равно минимальному весу ненулевого кодового слова.

Кодовое слово веса $1$ существовало бы тогда и только тогда, когда некоторый столбец $H$ равен нулю.
Но все столбцы матрицы $H$ ненулевые.
Значит, слов веса $1$ нет.

Кодовое слово веса $2$ существовало бы тогда и только тогда, когда сумма двух различных столбцов $H$ равна нулю.
Над $\F_2$ это означало бы, что эти два столбца равны.
Но все ненулевые столбцы различны.
Значит, слов веса $2$ нет.

Так как среди столбцов есть, например,
\[
    h_1,\quad h_2,\quad h_1+h_2,
\]
их сумма равна нулю:
\[
    h_1+h_2+(h_1+h_2)=0.
\]
Следовательно, существует кодовое слово веса $3$.

Итак,
\[
    d=3.
\]
\end{proof}

Поскольку
\[
    d=3=2\cdot 1+1,
\]
код Хэмминга исправляет одну ошибку.
Также он обнаруживает две ошибки, поскольку
\[
    d=3\ge 2+1.
\]

\section{Разреженность кодовых слов}

\subsection{Понятие разреженности}

В теории кодов с исправлением ошибок важна не только длина кодовых слов и их количество, но и то, насколько редко кодовые слова расположены во всём пространстве возможных слов.

Пусть задан $q$-ичный код
\[
    C\subseteq Q^n,
    \qquad |Q|=q.
\]
Всего в пространстве $Q^n$ существует
\[
    q^n
\]
слов длины $n$.
Если код содержит
\[
    M=|C|
\]
кодовых слов, то доля кодовых слов среди всех слов длины $n$ равна
\[
    \frac{M}{q^n}.
\]

\begin{definition}
\textbf{Плотностью кода} называется величина
\[
    \Delta(C)=\frac{|C|}{q^n}.
\]
\end{definition}

\begin{definition}
\textbf{Разреженностью кода} или \textbf{разреженностью множества кодовых слов} называется величина, обратная плотности:
\[
    S(C)=\frac{q^n}{|C|}.
\]
\end{definition}

Иными словами, разреженность показывает, сколько в среднем слов пространства $Q^n$ приходится на одно кодовое слово.

Если код имеет параметры
\[
    (n,M,d)_q,
\]
то
\[
    \Delta(C)=\frac{M}{q^n},
    \qquad
    S(C)=\frac{q^n}{M}.
\]

Если код линейный и имеет параметры
\[
    [n,k,d]_q,
\]
то
\[
    |C|=q^k.
\]
Поэтому
\[
    \Delta(C)=\frac{q^k}{q^n}=q^{k-n},
\]
а разреженность равна
\[
    S(C)=q^{n-k}.
\]

Число
\[
    n-k
\]
для линейного кода называется числом проверочных символов.
Таким образом, для линейного кода
\[
    S(C)=q^{n-k}.
\]

\begin{example}
Рассмотрим двоичный код Хэмминга
\[
    [7,4,3]_2.
\]
Он содержит
\[
    |C|=2^4=16
\]
кодовых слов.
Всего двоичных слов длины $7$:
\[
    2^7=128.
\]
Плотность кода:
\[
    \Delta(C)=\frac{16}{128}=\frac18.
\]
Разреженность:
\[
    S(C)=\frac{128}{16}=8.
\]
То есть в среднем на одно кодовое слово приходится $8$ слов пространства $\F_2^7$.
\end{example}

\subsection{Связь разреженности с исправлением ошибок}

Если код исправляет $t$ ошибок, то шары Хэмминга радиуса $t$ вокруг различных кодовых слов не пересекаются.

Объём шара Хэмминга радиуса $t$ в $q$-ичном пространстве длины $n$ равен
\[
    V_q(n,t)=\sum_{i=0}^{t}\binom{n}{i}(q-1)^i.
\]
Поэтому
\[
    |C|V_q(n,t)\le q^n.
\]
Отсюда
\[
    V_q(n,t)\le \frac{q^n}{|C|}=S(C).
\]

Следовательно, разреженность кода должна быть достаточно большой, чтобы вокруг каждого кодового слова можно было разместить шар исправляемых ошибок без пересечения с шарами других кодовых слов.

\begin{remark}
Чем больше исправляющая способность кода, тем больше должно быть минимальное расстояние между кодовыми словами и тем более разреженным становится множество кодовых слов.
\end{remark}

\subsection{Информационная скорость и разреженность}

Для кода с параметрами
\[
    (n,M,d)_q
\]
информационная скорость определяется как
\[
    R=\frac{\log_q M}{n}.
\]
Если код линейный $[n,k,d]_q$, то
\[
    R=\frac{k}{n}.
\]

Так как
\[
    M=q^{nR},
\]
то плотность кода равна
\[
    \Delta(C)=\frac{q^{nR}}{q^n}=q^{-n(1-R)}.
\]
Разреженность:
\[
    S(C)=q^{n(1-R)}.
\]

Таким образом:
\[
    R\text{ больше}
    \quad\Longleftrightarrow\quad
    \text{код плотнее},
\]
а
\[
    R\text{ меньше}
    \quad\Longleftrightarrow\quad
    \text{код разреженнее}.
\]

Высокая скорость означает, что код передаёт много полезной информации на один символ кодового слова, но обычно это снижает помехоустойчивость. Низкая скорость означает большую избыточность и большую потенциальную способность к исправлению ошибок.

\section{Границы для кодов}

\subsection{Граница Хэмминга}

Пусть $C$ --- $q$-ичный код длины $n$, исправляющий $t$ ошибок.
Объём шара Хэмминга радиуса $t$ равен
\[
    V_q(n,t)=\sum_{i=0}^{t}\binom{n}{i}(q-1)^i.
\]

\begin{theorem}[Граница Хэмминга]
Если существует $q$-ичный код с параметрами $(n,M,d)_q$, исправляющий $t$ ошибок, где
\[
    t=\left\lfloor\frac{d-1}{2}\right\rfloor,
\]
то
\[
    M\sum_{i=0}^{t}\binom{n}{i}(q-1)^i\le q^n.
\]
\end{theorem}

\begin{proof}
Шары Хэмминга радиуса $t$ вокруг различных кодовых слов не пересекаются, иначе одно и то же принятое слово могло бы быть декодировано двумя способами.

Каждый шар имеет объём
\[
    V_q(n,t)=\sum_{i=0}^{t}\binom{n}{i}(q-1)^i.
\]
Всего таких шаров $M$.
Все они лежат в пространстве $Q^n$, содержащем $q^n$ слов.
Следовательно,
\[
    M V_q(n,t)\le q^n.
\]
\end{proof}

\begin{definition}
Код называется \textbf{совершенным}, если в границе Хэмминга достигается равенство:
\[
    M\sum_{i=0}^{t}\binom{n}{i}(q-1)^i=q^n.
\]
\end{definition}

Коды Хэмминга являются совершенными одноошибочными кодами:
\[
    M=2^{2^r-r-1},\qquad n=2^r-1,\qquad t=1.
\]
Тогда
\[
    M(1+n)=2^{2^r-r-1}\cdot 2^r=2^{2^r-1}=2^n.
\]

\subsection{Граница Синглтона}

\begin{theorem}[Граница Синглтона]
Для любого $q$-ичного кода с параметрами $(n,M,d)_q$ выполнено
\[
    M\le q^{n-d+1}.
\]
Для линейного кода $[n,k,d]_q$ это эквивалентно
\[
    d\le n-k+1.
\]
\end{theorem}

\begin{proof}
Удалим из каждого кодового слова последние $d-1$ координат.
Останутся слова длины
\[
    n-d+1.
\]
Если два различных кодовых слова после удаления совпали, то они различались бы только в удалённых $d-1$ координатах.
Тогда расстояние между ними было бы не больше $d-1$, что противоречит минимальному расстоянию $d$.

Значит, после удаления координат все $M$ слов остаются различными.
Но слов длины $n-d+1$ над $q$-ичным алфавитом всего
\[
    q^{n-d+1}.
\]
Следовательно,
\[
    M\le q^{n-d+1}.
\]

Если код линейный, то
\[
    M=q^k.
\]
Тогда
\[
    q^k\le q^{n-d+1},
\]
откуда
\[
    k\le n-d+1,
\]
то есть
\[
    d\le n-k+1.
\]
\end{proof}

\section{Циклические коды}

\subsection{Определение}

\begin{definition}
Линейный код $C\le \F_q^n$ называется \textbf{циклическим}, если из
\[
    c=(c_0,c_1,\dots,c_{n-1})\in C
\]
следует
\[
    (c_{n-1},c_0,c_1,\dots,c_{n-2})\in C.
\]
\end{definition}

Слову
\[
    c=(c_0,c_1,\dots,c_{n-1})
\]
сопоставляется многочлен
\[
    c(x)=c_0+c_1x+\dots+c_{n-1}x^{n-1}.
\]
Циклический сдвиг соответствует умножению на $x$ по модулю
\[
    x^n-1.
\]

Поэтому циклический код можно рассматривать как идеал в кольце
\[
    \F_q[x]/(x^n-1).
\]

\begin{theorem}
Всякий циклический код $C$ является главным идеалом:
\[
    C=(g(x)),
\]
где $g(x)$ --- некоторый делитель многочлена $x^n-1$.
Многочлен $g(x)$ называется порождающим многочленом кода.
\end{theorem}

\begin{proof}
Кольцо $\F_q[x]$ является евклидовым, поэтому каждый его идеал главный.
Рассмотрим множество многочленов, соответствующих словам циклического кода, как идеал фактор-кольца
\[
    \F_q[x]/(x^n-1).
\]
Пусть $g(x)$ --- ненулевой многочлен наименьшей степени в этом идеале.
Докажем, что любой $c(x)\in C$ делится на $g(x)$ по модулю $x^n-1$.

Разделим $c(x)$ на $g(x)$:
\[
    c(x)=q(x)g(x)+r(x),
\]
где
\[
    \deg r<\deg g.
\]
Так как $c(x)$ и $q(x)g(x)$ лежат в идеале, то $r(x)$ также лежит в идеале.
Но $g(x)$ выбран ненулевым многочленом минимальной степени.
Следовательно,
\[
    r(x)=0.
\]
Значит, все слова кода порождаются $g(x)$.

Так как работа происходит по модулю $x^n-1$, должно выполняться
\[
    g(x)\mid x^n-1.
\]
\end{proof}

\section{Контрольные суммы и CRC}

\subsection{Контрольная сумма}

Контрольная сумма --- это короткая функция от сообщения, используемая для обнаружения ошибок.
Простейшая контрольная сумма может быть суммой байтов по модулю некоторого числа.

Такие методы обычно хорошо обнаруживают случайные искажения, но не всегда устойчивы к специально подобранным изменениям.

\subsection{CRC}

CRC, или cyclic redundancy check, основан на делении многочленов над $\F_2$.

Пусть сообщение представлено многочленом
\[
    m(x).
\]
Выбирается порождающий многочлен
\[
    g(x)
\]
степени $r$.
Передаваемое слово строится так:
\[
    x^r m(x)+r(x),
\]
где $r(x)$ --- остаток от деления $x^r m(x)$ на $g(x)$.
Тогда передаваемый многочлен делится на $g(x)$.

\begin{algorithm}[H]
\caption{CRC-кодирование}
\begin{algorithmic}[1]
\STATE Выбрать порождающий многочлен $g(x)$ степени $r$.
\STATE К сообщению приписать $r$ нулей, что соответствует умножению на $x^r$.
\STATE Разделить $x^r m(x)$ на $g(x)$ над $\F_2$.
\STATE Получить остаток $r(x)$ степени меньше $r$.
\STATE Заменить приписанные нули на биты остатка.
\end{algorithmic}
\end{algorithm}

\begin{example}
Пусть
\[
    m=1101,\qquad g=1011.
\]
Здесь
\[
    \deg g=3.
\]
Приписываем три нуля:
\[
    1101000.
\]
Делим $1101000$ на $1011$ по модулю $2$.
Остаток равен
\[
    001.
\]
Передаваемое слово:
\[
    1101001.
\]
При проверке получатель делит $1101001$ на $1011$.
Если остаток равен нулю, ошибка не обнаружена; если остаток ненулевой, ошибка обнаружена.
\end{example}

\section{Краткое сопоставление кодирования источника и канального кодирования}

\[
\begin{array}{p{5cm}|p{5cm}|p{5cm}}
\toprule
\textbf{Свойство} & \textbf{Кодирование источника} & \textbf{Канальное кодирование}\\
\midrule
Цель & Уменьшить объём данных & Защитить данные от ошибок\\
\midrule
Избыточность & Удаляется & Добавляется\\
\midrule
Типичные методы & Хаффман, арифметическое кодирование, LZ77, LZW, RLE & Коды Хэмминга, BCH, Рида--Соломона, CRC\\
\midrule
Ключевая величина & Энтропия & Минимальное расстояние\\
\midrule
Основное ограничение & $\overline l\ge H(X)$ & $d\ge 2t+1$ для исправления $t$ ошибок\\
\bottomrule
\end{array}
\]

\section{Минимальный набор формул}

\[
    H_D(X)=\sum_i p_i\log_D\frac1{p_i}.
\]

\[
    \overline l=\sum_i p_i l_i.
\]

\[
    \sum_i D^{-l_i}\le 1.
\]

\[
    H_D(X)\le \overline l < H_D(X)+1.
\]

\[
    d(x,y)=|\{i:x_i\ne y_i\}|.
\]

\[
    d(C)=\min_{x\ne y}d(x,y).
\]

\[
    \text{обнаружение }s\text{ ошибок}\iff d\ge s+1.
\]

\[
    \text{исправление }t\text{ ошибок}\iff d\ge 2t+1.
\]

\[
    V_q(n,t)=\sum_{i=0}^{t}\binom{n}{i}(q-1)^i.
\]

\[
    M V_q(n,t)\le q^n.
\]

\[
    d\le n-k+1.
\]

\[
    C=\{x\in\F_q^n:Hx^T=0\}.
\]

\[
    s=Hy^T=He^T.
\]

\end{document}
